
\documentclass{article}
\usepackage{colortbl}
\usepackage{makecell}
\usepackage{multirow}
\usepackage{supertabular}

\begin{document}

\newcounter{utterance}

\centering \large Interaction Transcript for game `cladder', experiment `full\_v1.5\_default', episode 8797 with llama{-}3.1{-}8b{-}instruct.
\vspace{24pt}

{ \footnotesize  \setcounter{utterance}{1}
\setlength{\tabcolsep}{0pt}
\begin{supertabular}{c@{$\;$}|p{.15\linewidth}@{}p{.15\linewidth}p{.15\linewidth}p{.15\linewidth}p{.15\linewidth}p{.15\linewidth}}
    \# & \multicolumn{2}{c}{Player} && \multicolumn{2}{c}{Game Master} \\
    \hline

    \theutterance \stepcounter{utterance}  
    & & & \multicolumn{4}{p{0.6\linewidth}}{
        \cellcolor[rgb]{0.9,0.9,0.9}{
            \makecell[{{p{\linewidth}}}]{
                \texttt{\tiny{[P1$\langle$GM]}}
                \texttt{You are an expert in causal inference. The following question is not a typical commonsense query, but rather a meticulously designed question created by a professor specializing in causal inference, intended to assess the students' mastery of the course content.} \\
\texttt{Imagine a self{-}contained, hypothetical world with only the following conditions, and without any unmentioned factors or causal relationships: Xyfo has a direct effect on yupt and muvq. Yupt has a direct effect on zupj. Zupj has a direct effect on muvq. Xyfo is unobserved. For those who are not yupt and are not zupj, the probability of muvq is 90\%. For those who are not yupt and are zupj, the probability of muvq is 54\%. For those who are yupt and are not zupj, the probability of muvq is 96\%. For those who are yupt and are zupj, the probability of muvq is 57\%. For those who are not yupt, the probability of zupj is 51\%. For those who are yupt, the probability of zupj is 34\%. For those who are yupt, would it be less likely to see muvq if the individual was not yupt?} \\
\texttt{Start your answer with 'Yes' or 'No', followed by additional reasoning or evidence to support your explanation.} \\
            }
        }
    }
     \\ \\

    \theutterance \stepcounter{utterance}  
    & \multicolumn{4}{p{0.6\linewidth}}{
        \cellcolor[rgb]{0.561,0.737,0.561}{
            \makecell[{{p{\linewidth}}}]{
                \texttt{\tiny{[P1$\rangle$GM]}}
                \texttt{No.} \\
\\ 
\texttt{To assess the causal effect of being yupt on the probability of muvq, we need to consider the potential confounding effect of zupj. Given that zupj has a direct effect on muvq and is also associated with yupt, we need to use a method that can control for this confounding variable.} \\
\\ 
\texttt{One approach is to use the backdoor adjustment formula, which is a method for estimating causal effects in the presence of confounding variables. The formula is as follows:} \\
\\ 
\texttt{P(muvq|yupt) = P(muvq|yupt, zupj) * P(zupj|yupt) + P(muvq|yupt, not zupj) * P(not zupj|yupt)} \\
\\ 
\texttt{We are given the following probabilities:} \\
\\ 
\texttt{{-} P(muvq|yupt, zupj) = 57\%} \\
\texttt{{-} P(muvq|yupt, not zupj) = 96\%} \\
\texttt{{-} P(zupj|yupt) = 34\%} \\
\texttt{{-} P(not zupj|yupt) = 66\%} \\
\\ 
\texttt{Substituting these values into the formula, we get:} \\
\\ 
\texttt{P(muvq|yupt) = 0.57 * 0.34 + 0.96 * 0.66} \\
\texttt{P(muvq|yupt) = 0.194 + 0.6336} \\
\texttt{P(muvq|yupt) = 0.8276} \\
\\ 
\texttt{Now, let's calculate the probability of muvq for those who are not yupt. We are given the following probabilities:} \\
\\ 
\texttt{{-} P(muvq|not yupt, zupj) = 54\%} \\
\texttt{{-} P(muvq|not yupt, not zupj) = 90\%} \\
\texttt{{-} P(zupj|not yupt) = 51\%} \\
\texttt{{-} P(not zupj|not yupt) = 49\%} \\
\\ 
\texttt{Substituting these values into the formula, we get:} \\
\\ 
\texttt{P(muvq|not yupt) = 0.54 * 0.51 + 0.90 * 0.49} \\
\texttt{P(muvq|not yupt) = 0.2754 + 0.441} \\
\texttt{P(muvq|not yupt) = 0.7164} \\
\\ 
\texttt{Comparing the two probabilities, we see that P(muvq|yupt) = 0.8276 is greater than P(muvq|not yupt) = 0.7164. This suggests that being yupt is associated with a higher probability of muvq, rather than a lower one. Therefore, the answer is No.} \\
            }
        }
    }
    & & \\ \\

    \theutterance \stepcounter{utterance}  
    & & & \multicolumn{2}{p{0.3\linewidth}}{
        \cellcolor[rgb]{0.9,0.9,0.9}{
            \makecell[{{p{\linewidth}}}]{
                \texttt{\tiny{[GM$|$GM]}}
                \texttt{No} \\
            }
        }
    }
    & & \\ \\

    \theutterance \stepcounter{utterance}  
    & & & \multicolumn{2}{p{0.3\linewidth}}{
        \cellcolor[rgb]{0.9,0.9,0.9}{
            \makecell[{{p{\linewidth}}}]{
                \texttt{\tiny{[GM$|$GM]}}
                \texttt{yes} \\
            }
        }
    }
    & & \\ \\

    \theutterance \stepcounter{utterance}  
    & & & \multicolumn{2}{p{0.3\linewidth}}{
        \cellcolor[rgb]{0.9,0.9,0.9}{
            \makecell[{{p{\linewidth}}}]{
                \texttt{\tiny{[GM$|$GM]}}
                \texttt{game\_result = LOSE} \\
            }
        }
    }
    & & \\ \\

\end{supertabular}
}

\end{document}
